% ---------------------------------------------------------------------
% Section 7: Conclusion
% ---------------------------------------------------------------------
\section{Conclusion}

This paper starts from an observation: dynamics-level control requires an acceleration-level task-space model, whereas dual-quaternion formulations typically stop at the velocity level. Filling in that order with third-order dual quaternions makes pose, velocity, and task-space acceleration available from a single kinematic product, lets a single multiplication generate the pose error and the twist error together, and removes the need for a separate acceleration-level error. The resulting structure endows the closed-loop storage function with an exact dissipation identity, so that nominal convergence, $L_2$ disturbance attenuation, and the steady bound for bias-type disturbances are proved in one framework and map onto one set of physical gains. Most valuable for tuning is the exactness of the attenuation metric: the certified $L_2$ gain equals the reciprocal of the smallest damping and coincides with the true $H_\infty$ norm of the linearized loop, so inferring a damping lower bound from a disturbance level is one division, not an estimate. The near-identity linearization also yields a practical warning---the rotational stiffness suffers a factor-$1/4$ reduction, which the hardware gain configuration compensates.

The hardware results support this reading. Under an equal gain budget the proposed law matches three representative baselines in tracking accuracy; the difference lies in provability: only the proposed law satisfies its attenuation certificate, damping condition, and sublevel margin in all 33 runs. The $\gamma$ ladder tightens the certified gain from $1/6\,\mathrm{s}$ to $1/15\,\mathrm{s}$ (boundary tier $1/18\,\mathrm{s}$), with the orientation transient improving monotonically. Meanwhile, the quasi-static residual is set by the balance between the firmware tether spring and friction and does not vary with the task stiffness---an experimental delineation of the applicability boundary of the static scaling law \eqref{eq:5.9}.

\textbf{Limitations and future work}: the strictification of the error dynamics (Appendix C.4) is incomplete and the stability conclusions are local; the hardware validation covers a single pose without exogenous disturbance injection, and quantitative on-hardware verification of the static scaling law requires a tether-free setup or a task stiffness well above the tether stiffness; a large-initial-error sweep for the geometric robustness is left as future work.
